\scanpage{15}

% Detached point at whole-source (199,104), below-left of the initial J;
% its textual status is uncertain, not its physical presence.
$\underset{\hspace{-0.6em}\unsure{\cdot|\phantom{\cdot}}}{\mathrm{J}}$e te signale que j'ai \unsure{réfléchi} aux Gr-champs sur $X$. Si $G$ est un Groupe sur $X$, $N$ un $G$-Module les Gr-champs sur $X$ ``épinglés par $G,N$'' forment a priori une $2$-\unsure{catégorie} et même une $2$-catégorie de Picard stricte, grâce à l'opération évident à la Baer. On trouve que le complexe \inserted{(de cochaines)} tronqué à l'échelon $2$ qui leur correspond est le tronqué
\[
\tau_{\leq 2}\bigl(\mathbb R\Gamma(B_{G}\mathbin{\mathrm{mod}}X,N)[1]\bigr)
\]
(\uline{NB} la cohomologie de $\mathbb R\Gamma(B_{G}\mathbin{\mathrm{mod}}X,N)$ commence en degré $1$) Plus \unsure{géométriquement}, un Gr-champ sur $X$ \unsure{épi-glé|épinglé} \inserted{par $(G,N)$} est essentiellement ``la même chose'' qu'une $2$-gerbe sur $B_{G}$, liée par $N$, et \deleted{\unsure{épi|pri}}\unsure{munie} d'une \unsure{trivialisation} au dessus de $X\approx B_{c}=(B_{G})_{/P}$ (où $P$ est le $\overset{\unsure{\deleted{\mathrm{s}}|\ast}}{\mathrm{o}}$bjet de $B_{G}$ ``torseur universel sous $G$''). Ces $2$-gerbes forment en fait une $3$-catégorie
% The uncertain crossed mark above the o of objet is at about (475,1362).
% Detached point at whole-source (175,1631), below de in the final line.
$\underset{\unsure{\cdot|\phantom{\cdot}}}{\text{de}}$ Picard a priori, mais il se
% archive sheet number: 14
